\documentclass[11pt]{article}
\usepackage[utf8]{inputenc}
\usepackage{amsmath, amssymb}
\usepackage{geometry}
\geometry{a4paper, margin=1in}
\usepackage{hyperref}

\title{Ullr's Secret Core Calculations Reference}
\author{}
\date{}

\begin{document}
\maketitle

\section*{Overview}
This document outlines the physical formulas and theoretical basis for the calculations implemented in \texttt{core.py}.

\section{Vapor Pressure and Humidity}
\subsection{\texttt{saturation\_vapor\_pressure(t\_celsius)}}
The saturation vapor pressure $e_s$ (in hPa) at a given temperature $T$ (in $^\circ\text{C}$) is calculated using the Magnus formula:
\[
e_s(T) = 6.112 \exp\left(\frac{17.67 T}{T + 243.5}\right)
\]

\subsection{\texttt{get\_dew\_point\_from\_rh(temp\_f, rh)}}
First, the actual vapor pressure $e$ is determined from the relative humidity ($RH$ in \%):
\[
e = e_s(T) \times \frac{RH}{100}
\]
The dew point temperature $T_d$ (in $^\circ\text{C}$) is derived by rearranging the Magnus formula:
\[
T_d = \frac{243.5 \ln(e / 6.112)}{17.67 - \ln(e / 6.112)}
\]

\subsection{\texttt{get\_rh\_from\_dew\_point(temp\_f, dew\_point\_f)}}
The relative humidity is calculated from the ratio of actual vapor pressure (saturation vapor pressure at the dew point $T_d$) to the saturation vapor pressure at air temperature $T$:
\[
RH = \left( \frac{e_s(T_d)}{e_s(T)} \right) \times 100\%
\]

\section{Wet Bulb Temperature}
\subsection{\texttt{psychrometric\_equation(tw\_c, t\_c, e, p\_hpa)} and \texttt{wet\_bulb\_f(...)}}
The wet bulb temperature $T_w$ is solved using the psychrometric equation:
\[
e_s(T_w) - P \cdot A \cdot (1 + 0.00115 T_w) \cdot (T - T_w) - e = 0
\]
where $P$ is the atmospheric pressure (hPa) and $A = 6.66 \times 10^{-4}$ is the psychrometric constant. The root-finding algorithm (\texttt{fsolve}) finds the $T_w$ that satisfies this equation.

\section{Atmospheric Pressure}
\subsection{\texttt{pressure\_at\_elevation(elevation\_ft)}}
Standard atmospheric pressure $P$ (hPa) at a given elevation $h$ (meters) is approximated by:
\[
P(h) = 1013.25 \left( 1 - 2.25577 \times 10^{-5} h \right)^{5.25588}
\]

\section{Solar Position}
\subsection{\texttt{estimate\_solar\_position(lat, lon, dt\_aware)}}
The solar elevation angle $\alpha$ and azimuth angle $\gamma_s$ are calculated from the current UTC time and location.
\begin{align*}
b &= \frac{2\pi}{365}(N - 81) \quad \text{($N$ is day of year)} \\
\delta &= 23.45^\circ \sin(b) \quad \text{(Declination)} \\
\text{EOT} &= 9.87 \sin(2b) - 7.53 \cos(b) - 1.5 \sin(b) \quad \text{(Equation of Time)} \\
\text{LST} &= \text{UTC Hour} + \frac{\text{Longitude}}{15} + \frac{\text{EOT}}{60} \quad \text{(Local Solar Time)} \\
\omega &= 15^\circ (\text{LST} - 12) \quad \text{(Hour Angle)}
\end{align*}
Elevation $\alpha$:
\[
\sin(\alpha) = \sin(\phi)\sin(\delta) + \cos(\phi)\cos(\delta)\cos(\omega)
\]
Azimuth $\gamma_s$:
\[
\cos(\gamma_s) = \frac{\sin(\delta) - \sin(\alpha)\sin(\phi)}{\cos(\alpha)\cos(\phi)}
\]

\section{Radiative Equivalent Temperatures}
\subsection{\texttt{calculate\_radiative\_equivalent\_temps(...)}}
To estimate the impact of solar and thermal radiation on snow surface temperatures, we compute equivalent temperature shifts $T_{SW,eq}$ (shortwave) and $T_{LW,eq}$ (longwave) using a heat transfer coefficient $K_{rad} = 8.0$.

\subsubsection*{Shortwave Equivalent}
The angle of incidence $\theta_i$ on a slope (angle $\beta$, aspect $\gamma_a$):
\[
\cos(\theta_i) = \cos(\beta)\sin(\alpha) + \sin(\beta)\cos(\alpha)\cos(\gamma_s - \gamma_a)
\]
Net shortwave radiation ($SW_{net}$) incorporating cloud cover fraction $c$ and surface albedo $\alpha_s$:
\[
SW_{net} = 1361 \cdot \cos(\theta_i) \cdot (1 - 0.65 c^2) \cdot (1 - \alpha_s)
\]
\[
T_{SW,eq} = \frac{SW_{net}}{K_{rad}}
\]

\subsubsection*{Longwave Equivalent}
Clear sky emissivity $\epsilon_{clear}$ depends on vapor pressure $e_a$:
\[
\epsilon_{clear} = 0.51 + 0.066 \sqrt{e_a}
\]
Total atmospheric emissivity $\epsilon_{atm} = (1 - c)\epsilon_{clear} + c$. 
The total incoming longwave radiation considers the Sky View Factor ($SVF = \cos^2(\beta / 2)$) and terrain emissivity ($\epsilon_{terrain} = 0.98$):
\[
LW_{net} = \sigma \left[ SVF \cdot \epsilon_{atm} T_{air}^4 + (1 - SVF) \cdot \epsilon_{terrain} T_{air}^4 - T_{snow}^4 \right]
\]
where $T_{snow} = 273.15$ K and $\sigma = 5.67 \times 10^{-8}$.
\[
T_{LW,eq} = \frac{LW_{net}}{K_{rad}}
\]

\section{Effective Temperature}
\subsection{\texttt{effective\_temperature\_f(...)}}
The final snow effective temperature combines the thermodynamic wet-bulb temperature with the radiative shifts:
\[
T_{eff} = T_w + T_{SW,eq} + T_{LW,eq}
\]

\section{Snow Consolidation Physics}
\subsection{Unit Dimensionality Conversion}
The integrals of effective temperature calculated in the model are natively in Fahrenheit-hours (F-hrs). However, standard hydrological and physical coefficients for snow melt and ice growth expect Degree-Days Celsius ($^\circ\text{C-days}$). The conversion factor is:
\[
1 \text{ } ^\circ\text{C-day} = 1.8 \times 24 = 43.2 \text{ F-hrs}
\]
Thus, an integral $I_{\text{F-hrs}}$ is converted to $I_{\text{C-days}}$ via:
\[
I_{\text{C-days}} = \frac{I_{\text{F-hrs}}}{43.2}
\]

\subsection{Physical Density and Coefficients}
Real snow density $\rho$ is calculated using Snow Water Equivalent (SWE, mm) and physical snow depth $h_0$ (cm):
\[
\rho = \frac{\text{SWE}/10}{h_0}
\]
The dynamic percolation coefficient $K_M$ (for melting) and freeze conductivity coefficient $K_F$ (for refreezing) are linearly scaled by the physical density:
\begin{align*}
K_M &= 2.0\rho + 0.1 \\
K_F &= 10.0\rho + 0.5
\end{align*}

\subsection{Melt and Refreeze Depths}
The physical melt depth $D_{melt}$ (in cm) follows the Degree-Day Method:
\[
D_{melt} = K_M \cdot I_{m,\text{C-days}}
\]
where $I_{m,\text{C-days}}$ is the melt integral ($T_{eff} > 32^\circ\text{F}$) converted to Degree-Days Celsius.

The physical refreeze depth $D_{freeze}$ (in cm) representing crust thickness follows a generalized Stefan Phase Change Equation:
\[
D_{freeze} = K_F \sqrt{I_{f,\text{C-days}}}
\]
where $I_{f,\text{C-days}}$ is the absolute freeze integral ($T_{eff} < 32^\circ\text{F}$) converted to Degree-Days Celsius.

\section{Dynamic Corn Window Model}
\subsection{Thermodynamic Capacity and Density}
The prime corn skiing window (represented by cumulative heat integrals in F-hrs) is dynamically modeled as a function of the physical snow density $\rho$.

\subsubsection*{Start Threshold (Cold Content Overcome)}
Denser snow has higher heat capacity and thermal conductivity, meaning it stores more overnight cold content. More energy is required to bring the surface layer to an isothermal melting state. The minimum required integral is modeled via linear interpolation based on two empirical boundaries:
\begin{itemize}
    \item $\rho = 0.35$ (typical spring snow) requires $60$ F-hrs.
    \item $\rho = 0.50$ (high-alpine firn) requires $80$ F-hrs.
\end{itemize}
This gives the linear relationship:
\[
I_{start}(\rho) = 60 + \frac{80 - 60}{0.50 - 0.35} (\rho - 0.35) = 133.33 \rho + 13.34
\]

\subsubsection*{End Threshold (Structural Integrity Limit)}
Denser snow has greater structural integrity and can tolerate a deeper melt layer without collapsing into unskiable slush. We define the maximum allowable melt depth $D_{max}$ (in cm) using linear interpolation based on the following empirical boundaries:
\begin{itemize}
    \item $\rho = 0.35$ supports a maximum melt depth of $1.66$ cm (corresponding to the standard $90$ F-hrs limit).
    \item $\rho = 0.50$ supports a maximum melt depth of $3.05$ cm (corresponding to an observed $120$ F-hrs limit).
\end{itemize}
The maximum allowable depth is:
\[
D_{max}(\rho) = 1.66 + \frac{3.05 - 1.66}{0.50 - 0.35} (\rho - 0.35) = 9.27 \rho - 1.58
\]
Setting the physical melt depth $D_{melt} = D_{max}(\rho)$ and rearranging the Degree-Day equation gives the maximum allowable heat integral:
\[
I_{end}(\rho) = \frac{D_{max}(\rho) \times 43.2}{K_M(\rho)} = \frac{(9.27 \rho - 1.58) \times 43.2}{2.0\rho + 0.1}
\]

\end{document}